\documentclass[../../../include/open-logic-section]{subfiles}

\begin{document}

\olfileid{sth}{replacement}{refproofs}
\olsection{પરિશિષ્ટ: પ્રતિસ્થાપન અંગેના પરિણામો}

આ વિભાગમાં $\ZF$માં પ્રતિબિંબન સાબિત કરીશું. પ્રતિબિંબન
કયા અર્થમાં પ્રતિસ્થાપનને સમકક્ષ છે તે પણ સાબિત કરીશું.
અને આ બધાંમાંથી પ્રતિબિંબન/પ્રતિસ્થાપનની શક્તિ વિશેનું
એક રસપ્રદ પરિણામ પણ સાબિત કરીશું. \emph{સાવધાન:
આ પાઠ્યપુસ્તકમાં આ સહેલાઈથી સૌથી અઘરું ગણિત છે.}

શરૂઆતમાં એક લેમા જોઈએ. ટૂંકાણ માટે તેમાં ચલો અથવા
વસ્તુઓના અનુક્રમ દર્શાવવા \emph{ઉપરરેખા}ના સંકેતનો
ઉપયોગ થાય છે. એટલે, ``$\overline{a}_k$'' એ
``$a_{k_1}$, \dots, $a_{k_n}$''નું ટૂંકું લખાણ છે,
જ્યાં $n$ સંદર્ભથી નક્કી થાય છે.

\begin{lem}\ollabel{lemreflection}
દરેક $1 \leq i \leq k$ માટે
$\phi_i(\overline{v}_i, x)$ સૂત્ર હોય. ત્યારે દરેક
$\alpha$ માટે એવો $\beta > \alpha$ છે કે કોઈ પણ
$\overline{a}_1, \ldots, \overline{a}_k \in V_\beta$
અને દરેક $1 \leq i \leq k$ માટે:
\[
	\exists x\phi_i(\overline{a}_i, x) \rightarrow (\exists x \in V_\beta) \phi_i(\overline{a}_i, x)
\]
\end{lem}

\begin{proof}
નીચે પ્રમાણે પદ $\mu$ વ્યાખ્યાયિત કરીએ:
$\mu(\overline{a}_1, \ldots,
\overline{a}_k)$ એ લઘુતમ તબક્કો $V$ છે જે દરેક
$1 \leq i \leq k$ માટે નીચેની શરત સંતોષે છે:
\[
\exists x\phi_i(\overline{a}_i, x) \rightarrow (\exists x \in V) \phi_i(\overline{a}_i, x)
\]
\sourcecorrection{OLMD-121}{મૂળ પ્રદર્શિત શરતમાં
છેલ્લે એક વધારાનો બંધ કૌંસ છે. અહીં અનુગામી
સૂત્રના પહેલેથી બંધ થયેલા કૌંસ પછીનું વધારાનું
ચિહ્ન દૂર કર્યું છે.}
બધાં $\overline{a}_1, \ldots, \overline{a}_k$
માટે $\mu(\overline{a}_1, \ldots, \overline{a}_k)$
અસ્તિત્વમાં છે તે સહેલાઈથી ચકાસી શકાય છે. હવે
પ્રતિસ્થાપન અને આપણા પુનરાવર્તન પ્રમેય વડે વ્યાખ્યા આપીએ:
\begin{align*}
	S_0 & = V_{\alpha+1}\\
	S_{n+1} & = S_n \cup \bigcup
	\Setabs{\mu(\overline{a}_1, \ldots, \overline{a}_k)}
	{\overline{a}_1, \ldots, \overline{a}_k \in S_n} \\
	S &= \bigcup_{m < \omega} S_m.
\end{align*}
\sourcecorrection{OLMD-122}{મૂળની છેલ્લી પંક્તિમાં
યોગગણની સીમા $m$ હોવા છતાં પદ $S_n$ હતું; ત્યાં
$n$ મુક્ત રહી જતો હતો. અહીં તે $S_m$ કર્યું છે.}
દરેક $S_n$, અને તેથી $S$ પોતે પણ, $V_\alpha$
પછીનો તબક્કો છે. હવે $\overline{a}_1$,
\dots,~$\overline{a}_k \in S$ નક્કી કરો; તેથી કોઈ
$n < \omega$ એવો છે કે $\overline{a}_1$,
\dots, $\overline{a}_k \in S_n$. કોઈ
$1 \leq i \leq k$ નક્કી કરો અને ધારો કે
$\exists x \phi_i(\overline{a}_i,x)$. રચના મુજબ
$(\exists x \in \mu(\overline{a}_1,
\ldots, \overline{a}_k))\phi_i(\overline{a}_i, x)$,
તેથી $(\exists x \in S_{n+1})\phi_i(\overline{a}_i, x)$,
અને તેથી $(\exists x \in S)\phi_i(\overline{a}_i, x)$.
આમ $S$ એ જોઈતું $V_\beta$ છે.
\end{proof}
\noindent
હવે \olref[sth][replacement][ref]{reflectionschema} સરળતાથી
સાબિત કરી શકીએ:

\begin{proof}[સાબિતી]
$\alpha$ નક્કી કરો. સામાન્યતા ગુમાવ્યા વિના આપણે
ધારી શકીએ કે $\phi$માં માત્ર $\exists$, $\lnot$
અને $\land$ તર્કસંયોજકો છે (કારણ કે આ અભિવ્યક્તિ
માટે પૂરતા છે). $\phi$ના બધા ઉપસૂત્રો જટિલતાના
ક્રમમાં $\psi_1, \ldots, \psi_k$ તરીકે ગણીએ,
જેથી $\psi_k = \phi$. \olref{lemreflection} મુજબ
એવો $\beta > \alpha$ છે કે કોઈ પણ
$\overline{a}_i \in V_\beta$ અને દરેક
$1 \leq i \leq k$ માટે:
\begin{align}\label{reflectionnicelybehaved}
	\exists x\psi_i(\overline{a}_i, x) \rightarrow
	(\exists x \in V_\beta) \psi_i(\overline{a}_i, x)\tag{*}
\end{align}
$\psi_i$ની જટિલતા પરના અનુમાનથી બતાવીશું કે કોઈ પણ
$\overline{a}_i \in V_\beta$ માટે
$\psi_i(\overline{a}_i) \leftrightarrow
\psi_i^{V_\beta}(\overline{a}_i)$. જો $\psi_i$
આણ્વિક હોય, તો આ સ્પષ્ટ છે. આ દ્વિદિશ શરતથી એ પણ
મળે છે કે જો $\psi_i$ આ ગુણ ધરાવતા ઉપસૂત્રોનો
નિષેધ કે સંયોજન હોય, તો $\psi_i$ને પણ આ ગુણ છે. તેથી
એકમાત્ર રસપ્રદ કિસ્સો પરિમાણીકરણનો છે. હવે
$\overline{a}_i \in V_\beta$ નક્કી કરો; ત્યારે:
\begin{align*}
	(\exists x \psi_i(\overline{a}_i, x))^{V_\beta}
	&\text{ iff }
	(\exists x \in V_\beta)\psi_i^{V_\beta}(\overline{a}_i, x)
	&&\text{by definition}\\
	&\text{ iff }
	(\exists x \in V_\beta)\psi_i(\overline{a}_i,  x)
	&&\text{by hypothesis}\\
	&\text{ iff }
	\exists x \psi_i(\overline{a}_i, x)
	&&\text{by \eqref{reflectionnicelybehaved}}
\end{align*}
આથી અનુમાન પૂરું થાય છે; $\psi_k = \phi$ હોવાથી
જોઈતું પરિણામ મળે છે.
\end{proof}

$\ZF$માં પ્રતિબિંબન સાબિત થયું. અમારી સાબિતી મુખ્યત્વે
\citet{Montague1961}ને અનુસરે છે. હવે $\Z$માં
પ્રતિબિંબનથી પ્રતિસ્થાપન નીકળે છે તે સાબિત કરવું છે.
સાબિતી સરળ ફેરફાર સાથે \citet{Levy1960}ને અનુસરે છે.

અહીં $\Z$માં કામ કરીએ છીએ, તેથી પ્રતિબિંબનને ઉપર
આપેલા બરાબર સ્વરૂપમાં રજૂ કરી શકતા નથી. કારણ કે
તે સ્વરૂપમાં ``$V_\alpha$''નો સંકેત વપરાયો છે, અને
$\Z$માં તેની વ્યાખ્યા આપી શકાતી નથી (જુઓ
\olref[sth][spine][zf]{sec}). તેથી તેના બદલે
પ્રતિબિંબનનું દેખીતી રીતે નબળું સ્વરૂપ નીચે મુજબ આપીએ:
\sourcecorrection{OLMD-123}{મૂળ અહીં પ્રતિસ્થાપનનું
નબળું સ્વરૂપ કહે છે, પરંતુ તરત આવતી વ્યાખ્યા નબળા
પ્રતિબિંબનની છે. અહીં તે વિષય-સૂચક શબ્દ સુધાર્યો છે.}

\begin{defish}
\emph{નબળું પ્રતિબિંબન.} કોઈ પણ સૂત્ર $\phi$ માટે
એવો પરંપરિત ગણ $S$ છે કે $0$, $1$ અને $\phi$ના
બધા પરિમાણો $S$ના !!{element}s છે, અને
$(\forall \overline{x} \in S)(\phi \liff
\phi^S)$.
\end{defish}

આમાંથી પ્રતિસ્થાપન સાબિત કરવા પહેલાં,
\citet[first
part of Theorem 2]{Levy1960}ને અનુસરીને બતાવીશું
કે બે સૂત્રોને એકસાથે ``પ્રતિબિંબિત'' કરી શકીએ છીએ:

\begin{lem}[$\Z + \text{Weak-Reflection}$માં.]\ollabel{lem:reflect}
કોઈ પણ સૂત્રો $\psi, \chi$ માટે એવો પરંપરિત
ગણ $S$ છે કે $0$ અને $1$ (અને સૂત્રોના કોઈ પણ
પરિમાણો) $S$ના !!{element}s છે, અને
$(\forall \overline{x} \in S)((\psi \liff \psi^S) \land (\chi
\liff \chi^S))$.
\end{lem}

\begin{proof}
$\phi$ને $(z = 0 \land \psi) \lor (z = 1 \land \chi)$
સૂત્ર તરીકે લો.

અહીં ટૂંકું લખાણ વાપર્યું છે; ``$z = 0$''ને
``$\forall t\, t \notin z$'' અને ``$z =1$''ને
``$\forall s(s \in z
\liff \forall t\, t \notin s)$'' તરીકે પૂરું લખવું
જોઈએ. પરંતુ $0, 1 \in S$ અને $S$ પરંપરિત હોવાથી,
આ સૂત્રો $S$ માટે \emph{નિરપેક્ષ} છે; એટલે કે
તેમના પરિમાણકો $S$ સુધી મર્યાદિત કરીએ કે નહીં,
તે જ વસ્તુ માટે તે સાચાં રહેશે.\footnote{વધુ ઔપચારિક
રીતે, આ બેમાંથી કોઈ પણ સૂત્ર $\xi$ હોય, તો
$\xi(z) \liff \xi^S(z)$.}

નબળા પ્રતિબિંબન મુજબ યોગ્ય $S$ એવો મળે છે કે:
\begin{align*}
	(\forall z, \overline{x} \in S)(&\phi \liff \phi^S)\\
	\text{i.e. }(\forall z, \overline{x} \in S)(&((z = 0 \land \psi) \lor (z = 1 \land \chi)) \liff {}\\
	&\phantom{(}((z = 0 \land \psi) \lor (z = 1 \land \chi))^S)\\
	\text{i.e. }(\forall z, \overline{x} \in S)(&((z = 0 \land \psi) \lor (z = 1 \land \chi))\liff {}\\
	&\phantom{(}((z = 0 \land \psi^S) \lor (z = 1 \land \chi^S)))\\
	\text{i.e. }(\forall \overline{x} \in S)(&(\psi \liff \psi^S) \land (\chi \liff \chi^S))
\end{align*}
બીજું વિધાન ત્રીજું વિધાન આપે છે, કારણ કે ``$z = 0$''
અને ``$z=1$'' $S$ માટે નિરપેક્ષ છે; ચોથું વિધાન
$0 \neq 1$ હોવાથી મળે છે.
\end{proof}\noindent હવે \citet[Theorem 6]{Levy1960}ની
સાબિતી અનુસરીને તેને સરળ બનાવતાં પ્રતિસ્થાપન મેળવી શકીએ:

\begin{thm}[$\Z$ અને નબળા પ્રતિબિંબન પર]\label{thm:replacement}
કોઈ પણ સૂત્ર $\phi(v,w)$ અને કોઈ પણ $A$ માટે,
જો $(\forall x \in A)\lexists![y][\phi(x,y)]$ હોય,
તો $\Setabs{y}{(\exists x \in A)\phi(x,y)}$ અસ્તિત્વમાં છે.
\sourcecorrection{OLMD-124}{મૂળના ગણ-નિર્માણ
સૂત્રમાં $\phi(x,y)$નું અંતિમ બંધ કૌંસ છૂટી ગયું છે;
અહીં તે ઉમેરી સૂત્ર પૂર્ણ કર્યું છે.}
\end{thm}

\begin{proof}
$(\forall x \in A)\lexists![y][\phi(x,y)]$ થાય એવો
$A$ નક્કી કરો અને નીચેના સૂત્રો વ્યાખ્યાયિત કરો:
\begin{align*}
	\psi &\text{ is } (\phi(x, z) \land A = A)\\
	\chi &\text{ is } \lexists[y][\phi(x, y)]
\end{align*}
\olref{lem:reflect} વાપરીએ. $A$ એ $\psi$નો
પરિમાણ હોવાથી, એવો પરંપરિત~$S$ મળે છે કે
$0, 1, A \in S$ (બીજા કોઈ પણ પરિમાણો સાથે), અને:
\[
	(\forall x,z \in S)((\psi \liff \psi^S) \land (\chi \liff \chi^S))
\]
ખાસ કરીને:
\begin{align*}
	(\forall  x, z \in S)(&\phi(x, z) \liff \phi^S(x, z))\\
	(\forall x \in S)(&\exists y\phi(x, y) \liff (\exists y \in S)\phi^S(x, y))
\end{align*}
આ બંનેને જોડીએ અને $A \in S$ તથા $S$ પરંપરિત
હોવાથી $A \subseteq S$ એ ધ્યાનમાં લઈએ:
\begin{align*}
	(\forall x \in A)(&\exists y\phi(x, y) \liff (\exists y \in S)\phi(x, y))
\end{align*}
હવે $(\forall x \in A)(\lexists![y \in S])\phi(x, y)$,
કારણ કે $(\forall x \in A)\lexists![y][\phi(x, y)]$.
પૃથક્કરણથી
$\Setabs{y \in S}{(\exists x \in A) \phi(x, y)} = \Setabs{y}{(\exists
x \in A) \phi(x, y)}$ મળે છે.
\end{proof}

\end{document}
